% R038：拓展阅读的偏导几何作用；等高线使用 c=a+xbar b 与 b 坐标，因而可轴对齐。
\begin{center}\begin{tikzpicture}[x=1cm,y=1.08cm,>=Stealth,every node/.style={font=\small}]
\path[use as bounding box] (0,0) rectangle (12.4,7.0);
\node[font=\large\bfseries,text=CProof] at (6.2,6.65) {拓展：两个偏导分别调截距与倾斜；误差碗只有一个最低点};
\draw[rounded corners=5pt,fill=purple!3,draw=CProof] (.15,.28) rectangle (6.02,6.12);
\node[font=\bfseries,text=CProof] at (3.08,5.72) {对 $a$：固定 $b$，直线上下平移};
\draw[->,black!65] (.70,1.02)--(5.45,1.02);\draw[->,black!65] (.70,1.02)--(.70,4.70);
\foreach \x/\y in {1.20/1.82,2.00/2.32,2.80/2.08,3.60/3.20,4.50/3.56}{\fill[blue!80!black] (\x,\y) circle (2.5pt);}
\draw[gray!65,thick] (.95,1.58)--(5.10,4.03);\draw[CStatement,very thick] (.95,1.35)--(5.10,3.80);\draw[gray!65,thick] (.95,1.12)--(5.10,3.57);
\fill[orange!90!black] (3.05,2.59) circle (3.4pt);\node[above left,text=orange!85!black] at (3.05,2.59) {中心};
\draw[red!78!black,thick] (2.0,2.32)--(2.0,1.97);\draw[red!78!black,thick] (3.6,3.20)--(3.6,2.92);\node[text=red!78!black,font=\scriptsize,anchor=west] at (3.16,3.38) {残差示意};
\draw[rounded corners=5pt,fill=purple!3,draw=CProof] (6.40,.28) rectangle (12.25,6.12);
\node[font=\bfseries,text=CProof] at (9.32,5.72) {$c=a+\bar x b$ 与 $b$ 的误差碗};
\draw[->,black!65] (6.98,1.05)--(11.68,1.05) node[right] {$c$};\draw[->,black!65] (6.98,1.05)--(6.98,4.82) node[above] {$b$};
\draw[CProof] (9.35,2.92) ellipse [x radius=.78,y radius=.42];\draw[CProof] (9.35,2.92) ellipse [x radius=1.42,y radius=.80];\draw[CProof] (9.35,2.92) ellipse [x radius=2.02,y radius=1.18];
\fill[orange!90!black] (9.35,2.92) circle (3.7pt);\node[above right,text=orange!85!black] at (9.35,2.92) {$(\bar y,\hat b)$};
\draw[->,red!75!black,very thick] (11.16,4.18)--(10.02,3.16) node[midway,above right,font=\scriptsize] {误差下降方向};\draw[->,red!75!black,very thick] (7.55,1.55)--(8.70,2.58);
\end{tikzpicture}\end{center}
